\section{Provenance of the spectral transport results}
\label{app:transport-provenance}

This appendix identifies the formal statements supporting the central
optimization theorems.  The claims in this section use three distinct
categories: established mathematical results used as inputs, new
statements proved in this project, and corrections of printed claims.
A Lean declaration verifies its formal proposition; correspondence to a
published statement is a separate source-review step.  The full
Davis--Kahan source audit appears in Appendix~\ref{app:dk-inventory}.

\paragraph{Classical geometry and optimality.}
Dixmier's relative-position analysis, Davis's principal-angle and
squared-displacement results, Halmos's two-subspace geometry, and
Kato's perturbation framework supply the established projection
machinery \citep{dixmier1948position,davis1958separation,halmos1969two,kato1966perturbation}.
Davis's 1958 Theorems~7.3--7.4 establish the Hilbert--Schmidt and
finite-dimensional Schatten-$p$ optimality cited in the main text.
Davis--Kahan (1970) supplies the operator perturbation theorems,
direct-rotation framework, and complex phase example
\citep{davis1970rotation}.  Bolotnikov--Li--Rodman (2001) poses the
ambient unitary minimal-distortion problem and proves Frobenius and
Q-norm results in its stated regimes \citep{bolotnikov2001minimal}.
The project does not assign originality to those results.

\paragraph{Closed-family finite-dimensional value and phase threshold.}
Theorem~\ref{thm:family-finite} corresponds to
\path{SmoothedPowerTransport.exists_unitary_minimizing_transport_principalAngles_of_nonneg}
in
\path{DavisKahan/FiniteDimensional/DirectRotation/SmoothedPowerBoundary.lean}.
Its positive-smoothing structural argument is formalized by
\path{SmoothedPowerTransport.spectral_decomposition_of_isLocalMinOn_transports}
in
\path{DavisKahan/FiniteDimensional/DirectRotation/SmoothedPowerSpectralDecomposition.lean}.
Theorem~\ref{thm:family-canonical} is
\path{SmoothedPowerTransport.nonacuteDirectRotation_isMin_iff_principalAngles_critical}
in
\path{DavisKahan/FiniteDimensional/DirectRotation/SmoothedPowerCanonicalCriterion.lean}.
The scalar one-crossing analysis uses
\path{DavisKahan/FiniteDimensional/DirectRotation/SmoothedPowerPhaseSelection.lean}
and the general calculus lemma in
\path{ForTauCeti/Analysis/Calculus/SingleCrossingMinimum.lean}.
These are additional transport theorems; neither Davis (1958) nor
Bolotnikov--Li--Rodman (2001) is credited with this particular
closed-family formula.

\paragraph{Infinite-dimensional ideal and minimum.}
Theorem~\ref{thm:family-infinite} is supported by
\path{SmoothedPowerTransport.exists_unitary_minimizing_transport_of_projection_cost_ne_top}
in
\path{DavisKahan/InfiniteDimensional/DirectRotation/SmoothedPowerGlobalAttainment.lean}.
The finite-cost equivalence is
\path{SmoothedPowerTransport.finite_minimum_iff_projection_cost_and_defect_balance}
in the same infinite-dimensional development.  The finite-section
lower-bound principle is
\path{TauCeti.tsum_angle_penalty_le_displacement_of_finite_transport_bound}
in
\path{ForTauCeti/Analysis/InnerProductSpace/Projection/SpectralTransportLimit.lean}.
The relative-index and defect obstruction is credited to the classical
literature \citep{avron1994index}; the attained angle-series value and
the change of ideal with smoothing are results of this development.

\paragraph{Real matching and signed-root compatibility.}
The verified unrestricted real Schatten matching result is
\path{RealTransport.realTransportPowerMinimum_eq_matching}
in
\path{DavisKahan/FiniteDimensional/DirectRotation/RealNuclearMatching.lean},
with explicit ambient attainment proved in the companion real-matching
modules.  The exact nuclear real--complex parity statement is
\path{RealTransport.realTransportPowerMinimum_one_eq_sine_sum_iff_even_multiplicities}
in
\path{DavisKahan/FiniteDimensional/DirectRotation/RealNuclearFieldEquality.lean}.
The arbitrary-dimensional real nuclear equality criterion is
\path{exists_real_nuclear_transport_eq_projection_sub_iff_even_source_multiplicity}
in
\path{DavisKahan/Geometry/Polar/NuclearRealRootCriteria.lean}.
The signed-root equality-attainer correspondence is
\path{nuclearSignedSineEquivProjectionBoundAttainers} in the polar
subsystem.  These results use classical principal-angle and complex
structure theory; no earlier unrestricted real matching or nuclear
parity theorem was identified in the sources examined for this paper.
That is a bounded historical assessment, rather than a proof of first
priority.

\paragraph{Specialized optimizer classifications.}
\Cref{sec:schatten-optimization,sec:smoothed-transport} contain the
full classification and local landscape theorems in their established
subfamilies, with proof modules under
\path{DavisKahan/FiniteDimensional/DirectRotation/} and
\path{DavisKahan/InfiniteDimensional/DirectRotation/}.
These stronger descriptions have not been formally generalized to
all $1\le p\le2$, $\tau\ge0$.  The unverified generic-penalty
fragment \path{papers/dk_formalization_journal/spectral_penalty_transport.tex}
is excluded from this manuscript and contributes no claim.

\paragraph{Formal ancestry and scope.}
The development uses Mathlib's Hilbert and continuous-functional-calculus
interfaces \citep{mathlib2020,dedecker2025cfc}; parts of its spectral
measure and unbounded self-adjoint infrastructure were adapted from
Spectra \citep{bornemann2026spectra}.  Proof compilation and the
separate comparison of Lean statements with published hypotheses are
documented in Appendix~\ref{app:formalization-process}.
